\section*{Abstract}

This diploma thesis presents the design, implementation, and evaluation of an \ORMName{} library for C++23~\cite{cpp23standard}. The work is organised around two complementary axes. The first is compile-time modelling of data and queries, in which the logical structure of data access is not described through strings interpreted at execution time, but through statically typed \code{constexpr} constructions. The second is a coroutine-based asynchronous architecture through which the same type-constrained model is extended toward non-blocking or pseudo-non-blocking execution according to the capabilities of the concrete backend driver.

The developed library uses a static data model expressed in C++ through constructs such as \code{property<T, Name>}, \code{relationship<...>}, and \code{table\_name\_trait<T>}, together with a typed intermediate representation of queries realised through \code{select\_query}, \code{insert\_query}, \code{update\_query}, and \code{delete\_query}. This logical layer is independent of the concrete database type. Materialisation toward relational and non-relational systems is performed by specialised \code{connector\_trait<DB>} implementations, which declare supported capabilities and translate the same logical query into Structured Query Language (SQL), Binary JSON (BSON)-like filters, Redis commands, Cypher, or Cassandra Query Language (CQL) depending on the selected backend.

A significant extension of the architecture is its asynchronous layer, built on coroutines, \code{Task<T>}, \code{ThreadPool}, \code{async\_db<DB>}, and coroutine-aware connection management. The thesis shows that a single asynchronous abstraction can be realised in two ways: by offloading blocking operations to a thread pool and by using native non-blocking interfaces for drivers that provide them. In this way, the work does not only describe a logical ORM abstraction, but also studies the execution model through which that abstraction remains practically useful under concurrent workloads.

The thesis analyses both relational scenarios for SQLite, PostgreSQL, and MySQL, and non-relational scenarios for MongoDB, Redis, Neo4j, and Cassandra. It demonstrates how the same C++ model dictates the logical schema while being materialised differently as a table, document, key-value structure, graph node, or wide-column row. It also examines the native client libraries and their execution semantics in order to determine where the library can honestly offer native asynchronous behaviour and where the correct choice remains a synchronous driver with controlled offload execution.

Validation of the work is grounded in the repository's unit and integration tests, while the empirical evaluation is extended through benchmark scenarios comparing synchronous and asynchronous execution. These cover query construction, capability-based restriction of operations, prepared queries, migrations, thread safety, coroutine infrastructure, and the behaviour of concrete connectors with either live or mock drivers. This makes it possible to formulate not only an architectural model, but also concrete criteria for a “fully implemented and functioning” connector and for a practically justified asynchronous layer.

The main contribution of the thesis lies in four directions. First, it proposes a practically usable model for an \ORMName{} library that combines static type safety with a backend-independent architecture. Second, it shows how a C++ data model can serve as the primary source of a logical schema across different classes of databases. Third, it develops a coroutine-based asynchronous architecture that extends the connector model without compromising compile-time constraints. Fourth, it formalises a connector contract through which the library can be extended with new backend implementations and asynchronous capabilities without changing the user-facing query API.

\vspace{1cm}

\textbf{Keywords:} C++23, \ORMName{}, static type safety, coroutines, asynchronous execution, \code{connector\_trait}, relational and non-relational databases, empirical evaluation
